\chapter{Considerações Finais}
\label{cap:conclusao}

Este trabalho apresentou a \textit{PolyglotImportCSV}, uma ferramenta de importação
declarativa que materializa um mesmo conjunto de dados em formato CSV em cinco SGBDs de
paradigmas distintos --- PostgreSQL, MongoDB, Cassandra, Redis e Neo4j --- a partir de dois
arquivos de configuração validados estaticamente. Partindo da fundamentação sobre
persistência poliglota e do estado da arte em importação e modelagem de dados heterogêneos,
o trabalho evoluiu da proposta a uma implementação funcional, avaliada empiricamente.

As principais contribuições são: (i) um formato de configuração JSON, separado em mapeamento
e conexão e validado por JSON Schema antes de qualquer escrita, com dois modos de entrada ---
um arquivo por entidade (multifonte) e um único CSV combinado; (ii) o mapeamento declarativo
de colunas com conversão nativa aos tipos de cada SGBD; (iii) uma camada de observabilidade e
um subsistema de métricas que instrumentam cada fase da importação; (iv) dois eixos de
desempenho --- a estratégia de escrita otimizada, com vetorização e gravação em lote, e o
modelo de execução em fluxo, de memória limitada; e (v) uma avaliação empírica com dados reais
que quantifica o custo de importar para os cinco paradigmas e o ganho de cada eixo.

A avaliação (Capítulo~\ref{cap:avaliacao}) confirmou os objetivos traçados. A estratégia
otimizada acelerou a importação em cerca de nove vezes sobre a linha de base ingênua, e o modo
de execução em fluxo reduziu o pico de memória em mais de doze vezes no maior volume medido,
mantendo-o constante em relação ao tamanho do arquivo --- o que habilita a importação de
arquivos maiores que a memória disponível. A equivalência dos dados gravados entre os dois
modos de importação e entre os dois modelos de execução, assegurada por testes automatizados,
permitiu tratar a escolha entre eles como uma decisão de custo, não de correção.

As limitações reconhecidas dizem respeito sobretudo ao ambiente de avaliação: os experimentos
correram em uma única máquina, com os cinco SGBDs em contêineres partilhando os mesmos
recursos, o que penalizou em especial a escrita no Cassandra --- uma restrição do ambiente, e
não uma propriedade da ferramenta. Além disso, a suíte de testes permanece majoritariamente
unitária, sem integração contínua contra os SGBDs conteinerizados.

Como trabalho futuro (Capítulo~\ref{cap:futuros}), destacam-se uma interface gráfica que torne
a ferramenta acessível a usuários menos familiarizados com a linha de comando, o aprofundamento
da avaliação em volumes maiores e em ambientes com SGBDs dedicados, e a generalização do núcleo
de importação a novos destinos. Consolida-se, assim, a \textit{PolyglotImportCSV} como artefato
de pesquisa e de código aberto alinhado à literatura sobre persistência poliglota e modelagem
de dados heterogêneos.
